\section{One physical profile for the six original preparations}
Fix a common model profile before evaluating any lane:
\[
 \begin{gathered}
 N_\alpha(\mathbf x)=1-\alpha V_{\mathbf e}(\mathbf x),\qquad
 V_{\mathbf e}=\frac{\sqrt2d}{\sqrt{d^2+|\mathbf x-d\mathbf e|^2}},\\
 \alpha=10^{-3},\quad d=10^{30}R,\quad |\mathbf e|=1,
 \quad N_\alpha(0)=n=0.999.
 \end{gathered}
\]
This softened Plummer profile is smooth and globally positive:
$1-\sqrt2\alpha\le N_\alpha<1$. Its positive Laplacian is consistent with
a finite positive-density body under a declared weak-potential Poisson
interpretation. That interpretation is an additional receiving assignment,
not a gravity law derived from the original source ledger.

The original radius, propagation speed, action chart and clock window
$T_0=573049\pi$ remain common to every preparation. For a fully specified
shared SI fixture, choose $R_{\rm phys}=10^{-19}\,\mathrm m$,
$c_{\rm phys}=299792458\,\mathrm{m\,s^{-1}}$, and explicitly assign
$I_u=\hbar$, with $h=6.62607015\,10^{-34}\,\mathrm{J\,s}$. Then
\[
 \ell_u=R_{\rm phys}/R,\qquad t_u=\ell_u/c_{\rm phys},
 \qquad E_u=\hbar/t_u.
\]
The resulting displays are $t_u\simeq2.03959635\,10^{-27}\,\mathrm s$,
$E_u\simeq5.17049276\,10^{-8}\,\mathrm J$,
$\epsilon E_u\simeq9.02277600\,10^{-14}\,\mathrm J$,
$T_0t_u\simeq3.67185784\,10^{-21}\,\mathrm s$, and
$d_{\rm phys}=10^{11}\,\mathrm m$. Outward intervals travel with the
source bundle. This is one declared calibration fixture shared by all six
lanes; the microscopic assignment of the original action and ruler remains
a separate source question.

At the source center,
\[
 \nabla N_\alpha=-\frac{\alpha\mathbf e}{2d},\qquad
 \operatorname{Hess}N_\alpha=\frac{\alpha}{2d^2}
       (I-\tfrac32\mathbf e\mathbf e^{\rm T}).
\]
The gradient has zero first compression but retains its full local angular
complement. For $\beta=\angle(\mathbf e,\mathbf e_z)$ the quadratic
compression is
\[
 A(r)-nI=\frac{\alpha r^2}{56d^2}
       (9-6\cos^2\beta)I+O(r^4/d^4).
\]
Parallel, transverse and $\cos^2\beta=1/3$ controls have coefficients
$3/56,9/56,7/56$ times $\alpha r^2/d^2$. The last angle cancels the
trace-free tidal part, while the softened body's positive density leaves
an isotropic term.

A true leading tidal null is supplied by a second bounded physical
control: a uniform body centered at $d\mathbf e$ with radius $a=d/2$.
Writing $s=|\mathbf x-d\mathbf e|$,
\[
 N_\alpha=\begin{cases}
 1-\alpha d/s,&s\ge a,\\
 1-\alpha d(3a^2-s^2)/(2a^3),&s<a.
 \end{cases}
\]
It is globally positive and $C^1$, with minimum $1-3\alpha$ and continuous
normal derivative at its surface. The clock lies outside the body. Here
the quadratic coefficient is
$\alpha r^2(1-3\cos^2\beta)/(7d^2)$, with a genuine leading null at
$\cos^2\beta=1/3$. The finite globally bounded interior is retained.

\section{An evaluated first response with the interacting transient}
For the positive Bessel/TT receiver, the complete form is exactly
\[
 H_\alpha=(1-\alpha)H_0+\alpha r_B,\qquad
 r_B[\Phi,\Psi]=\langle T\Phi,J^*(1-V)JT\Psi\rangle.
\]
The original even source energy, both interactions and both counterterms
remain. Equality to a complete canonical local quantum stress is not assumed.
Let $\chi(t)=U_0(t)\Psi_{\rm prep}$ and let $z_1$ denote the derivative
of its actual current with respect to $\alpha$ at zero. Then
\[
 z_1(T)=-T\dot z_0(T)+z_{1,\rm inh}(T).
\]
The first term retains the exact interacting transient. The inhomogeneous
term is the weak two-time response pairing with $r_B$.

Choose the single common near radius $L=10^{18}R_-$, where $R_-$ is the
original radius's lower endpoint. Taylor's theorem and gradient cancellation
give the complete profile bounds
\[
 b_{\rm near}=\frac{3}{2\,10^{24}},\qquad b_{\rm far}=1.
\]
They are bounds for the derivative profile $1-V$, so the finite perturbation
pays $\delta_{\rm near}=\alpha b_{\rm near}$ and
$\delta_{\rm far}=\alpha b_{\rm far}$. The exterior is retained.

The original certificate gives $K[\chi(t)]\le\kappa_0$ and
$K_{\rm out}[\chi(t)]\le\eta_{\rm ref}(L)$ throughout the window.
For $B_T(s)=U_0(T-s)^*ZU_0(T-s)$, the current commutator and energy
conservation give, for $B_T$ and its adjoint,
\[
 K[B_T(s)\chi(s)]\le k_B
  =(2\sqrt{\kappa_0}+2\sqrt c)^2+8\epsilon.
\]
Each first-variation pairing has one original trajectory factor. Apply
its actual tail bound to that factor and $K_{\rm out}\le K\le k_B$ to
the current-inserted factor. This avoids assuming an uncomputed tail for
the inserted state and gives
\[
 \begin{aligned}
 |z_{1,\rm inh}(T)|&\le D_1,\\
 D_1&=2T\sqrt{k_B}
 [b_{\rm near}\sqrt{\kappa_0}
          +b_{\rm far}\sqrt{\eta_{\rm ref}(L)}].
 \end{aligned}
\]
With $m=M_--\mathcal E(T)>0$, the actual first phase derivative is enclosed by
\[
 \theta_1(T)=-T\nu_0(T)+[-D_1/m,D_1/m].
\]
Every quantity is evaluated with outward rational bounds from the original
six preparations. The following displays round outward:
\begin{center}\small
\begin{tabular}{c rr}\toprule
Lane&Lower $\theta_1$&Upper $\theta_1$\\\midrule
01&$-6.283185603$&$-6.283184981$\\
02&$-6.283185904$&$-6.283184578$\\
03&$-6.283197797$&$-6.283152294$\\
12&$-6.283185624$&$-6.283184955$\\
13&$-6.283196966$&$-6.283154685$\\
23&$-6.283196110$&$-6.283157146$\\\bottomrule
\end{tabular}
\end{center}
The complete complex $z_1$ is also enclosed: use
$|\dot z_0-\dot h|\le\Delta\mathcal E+2|g|W_D$ and the exact hard
harmonic $h=M e^{-\ii\Omega T}$. Its endpoint arc has angle
$T(\Delta-\Omega)$, so elementary sine/cosine bounds and the disk
$T(\Delta\mathcal E+2|g|W_D)+D_1$ give the retained complex rectangles.
All six imaginary parts of $z_1$ are strictly positive; its real part is
strictly negative in lanes03,13,23 and unresolved in the weaker lanes.
The inhomogeneous correction's own sign and its cross-lane differences
remain unresolved; these signed enclosures do not identify that correction
with its uniform contribution.

\section{Finite profile responses and a complete error budget}
The same actual source preparation evolves under the finite profile
$\alpha=10^{-3}$. Compare it with $\chi(nt)$, $n=.999$, and write
$e(t)=\|\Psi_\alpha(t)-\chi(nt)\|$. Static energy conservation and the
weak overlap identity give a strict sufficient first-crossing condition:
\[
 \begin{aligned}
 e_*^4>{}&4T^2 b_\alpha^2(3\kappa_0+4\epsilon e_*),\\
 b_\alpha={}&\delta_{\rm near}\sqrt{\kappa_0}
                 +\delta_{\rm far}\sqrt{\eta_{\rm ref}(L)}.
 \end{aligned}
\]
Its verified positive margins imply $e(t)<e_*$ throughout the window.
The current error is at most $4e_*$ and the phase error at most
$\arcsin(4e_*/m)$. The certificate uses the rational majorant
$(4e_*/m)/(1-4e_*/m)$ and verifies nonvanishing current at every prefix.
The source/reference/window part is retained exactly:
\[
 \theta_0(nT)-\theta_0(T)=-\int_{nT}^{T}\nu_0(t)\,\dd t.
\]
Thus the full finite phase responses, including the physical-profile
remainder, are enclosed as follows:
\begin{center}\small
\begin{tabular}{c rr}\toprule
Lane&Lower $\Delta\theta$&Upper $\Delta\theta$\\\midrule
01&$-.006285093$&$-.006281277$\\
02&$-.006287000$&$-.006279370$\\
03&$-.006313697$&$-.006252653$\\
12&$-.006285093$&$-.006281277$\\
13&$-.006298438$&$-.006267914$\\
23&$-.006298439$&$-.006267915$\\\bottomrule
\end{tabular}
\end{center}
All are strictly negative. The profile comparison errors are at most
$1.908,3.815,30.52,1.908,15.26,15.26$ microradians in lane order.
These are actual-state receiving-form certificates, not achieved apparatus
precision. The compact-body control also passes all six sign tests with
$b_{\rm near}=2\,10^{-24}$ and $b_{\rm far}=2$.

An explicit finite phase linearization payment is
\[
 \begin{aligned}
 |\Delta\theta_\alpha-\alpha\theta_1|\le{}&
  \phi_{\rm prof}+\alpha D_1/m\\
  &+\alpha T(\nu_+-\nu_-).
 \end{aligned}
\]
A complex-current finite remainder is retained separately in the source
bundle. A universal-coupling null must compare
$\theta_{\alpha,l}(T)$ with $\theta_{0,l}(nT)$, preserving each lane's
interacting transient. Pairwise differences of those residuals pay the sum
of the two profile errors together with preparation, timing, readout and
shared-reference covariance. At first order the corresponding residual is
$\theta_{1,l}+T\nu_{0,l}(T)$, bounded by $D_{1,l}/m_l$.

For the explicitly chosen receiving form, \textbf{the test result suggests
strong contact with the concept}: both the actual first response and a
finite perturbation have a certified sign in all six lanes under one
positive physical profile. The unresolved local quantum-stress terms and
the inhomogeneous susceptibility remain separately named recipients.
